\newif\ifglobal

%%%%%%%%%%%%%%%%%%%%%%%
%%% Compile to view %%%
%%%%%%%%%%%%%%%%%%%%%%%

%%%%%%%%%%%%%%%%%%%%%%%%%%%%%%%%%%%%%%%%%%%%%%%%%%%
%%% Readme:                                     %%%
%%% https://www.overleaf.com/read/bwdjvfjksvjy/ %%%
%%%%%%%%%%%%%%%%%%%%%%%%%%%%%%%%%%%%%%%%%%%%%%%%%%%

% >>> M?: marks a module setting number or important notice

% >>> M4: Set {./relative-path} to external files for this module <<<
\def \modulefiles {./11-INTMTRX}
% To add an external file, use:
% (Replace '---' with actual filename)
% '\input{\modulefiles/tables/---.tex}' for tables
% '\includegraphics{\modulefiles/figures/---}' for figures
% '\subfile{\modulefiles/---}' for register files

% >>> M1: This module is in global project? <<<
% 'YES' - uncomment; 'NO' - comment out
\globaltrue

\ifglobal
    \documentclass[main\_\_EN.tex]{subfiles}

\else
    \documentclass[openany, 10pt]{book}

    % >>> M2: In ./00-shared/preamble-trm-module, also find
    %         settings for visibility of todo notes, watermark <<<
    \usepackage{preamble-trm-module}

    % >>> M3: Temporary labels in English,
    %         you can add your own labels there <<<
    \myexternaldocument{./00-shared/temp-labels-en}

\fi %\ifglobal


\setcounter{secnumdepth}{4}


% Define formatting of parts and chapters in text
\titleformat{\chapter}
  [display]
  {\LARGE\bfseries}
  {\color{AllportsColor}Chapter \thechapter}
  {1em}
  {\color{AllportsColor}}
\titlespacing*{\chapter}{0pt}{0pt}{20pt}

% Define headers
\usepackage{fancyhdr}
\usepackage{lipsum}
\renewcommand{\chaptermark}[1]{\markboth{Chapter\, \thechapter\, \ #1}{}}
\fancypagestyle{plain}{%
  \fancyhead{}%
  %\fancyhead[R]{\Navigation}
  \fancyhead[L]{%
    \nouppercase{%
      \ %
      \normalsize\leftmark
    }%
  }%
}
\pagestyle{plain}


% >>> M5: If this module needs any updates in
% './00-shared/preamble-trm-module.sty'
% (include additional package, etc.), contact your module PM
% or create the file './00-shared/changelog.tex' make the
% necessary updates yourself and document them in this file <<<

% Make text left-aligned and allow latex to break pages naturally, comment out for full justification
\raggedright
\raggedbottom


\begin{document}


% Sets language into which pre-defined bits of text are translated
\selectlanguage{english}

% Do not delete this block
\ifglobal
% Add the following front matter if
% this module is outside of global project
\else
    \InputIfFileExists{readme}{}{}
    \listoftodos
    {\let\clearpage\relax \listofdonetodos}
    \tableofcontents
    %\listoftables
    %\listoffigures
\fi









%%%%%%%%%%%%%%%%%%%%%%%%%
%%% TRM Module Begins %%%
%%%%%%%%%%%%%%%%%%%%%%%%%

\hypertarget{intmtrx}{}
\chapter{Interrupt Matrix}
\label{mod:intmtrx}
\section{Overview}
The interrupt matrix embedded in \chipname{} routes one or multiple peripheral interrupt sources to any of the ESP-RISC-V CPU's peripheral interrupts.

The \chipname{} has 84 peripheral interrupt sources that can be routed to any of the 32 CPU peripheral interrupts using the interrupt matrix.

\begin{tiplisting}
This chapter focuses on how to map peripheral interrupt sources to the CPU interrupts. For information about interrupt configuration, vector, and interrupt handling operations recommended by the ISA, please refer to Chapter \ref{mod:riscvcpu} \textit{\nameref{mod:riscvcpu}}  > Section \ref{sec:riscvcpu-int-ctrl} \textit{\nameref{sec:riscvcpu-int-ctrl}}.
\end{tiplisting}
\vspace{-1em}

\section{Terminology}\label{sec:intmtx-int-terminology}

The following terms related to interrupts are defined in the context of the \textit{\chipname{} Technical Reference Manual} to help readers better understand this document:

\subsection{Interrupt}

An interrupt refers to the event or condition that occurs, causing the CPU to temporarily suspend its current execution and handle a higher-priority task. It is a mechanism that allows the CPU to respond to specific events promptly.

The \textit{\chipname{} Technical Reference Manual} may use the term ``interrupt" in a broader sense to refer to both interrupt signal and interrupt source.

\subsection{Interrupt Signal/Interrupt Source}

Interrupt signal and interrupt source are only defined from different perspectives and mean the same thing.

From the perspective of peripherals, interrupt signals are generated by the peripheral's internal interrupt sources and are sent to the interrupt matrix.

From the perspective of the interrupt matrix, it receives the interrupt signals sent from the peripheral and considers them interrupt sources. The interrupt matrix then outputs CPU peripheral interrupt signals to the CPU.

From the perspective of the CPU, the interrupt signals from the interrupt matrix become sources and are sent to the CPU core together with the core local interrupt sources.

\subsection{Interrupt Flow in \chipname{}}

Figure \ref{fig:interrupt-term} shows the interrupt flow in \chipname{}.

\begin{figure}[H]
    \centering
    \includegraphics[width=1\textwidth]{\modulefiles/figures/ESP32-H2-interrupt-terms.eps}
    \caption{Interrupt Flow in \chipname{}}
    \label{fig:interrupt-term}
\end{figure}

\section{Features}

The interrupt matrix embedded in the \chipname{} has the following features:

\begin{itemize}
    \item 84 peripheral interrupt sources accepted as input
    \item 32 HP CPU peripheral interrupts generated to the HP CPU as output
    \item Current interrupt status query of peripheral interrupt sources
    \item Multiple interrupt sources mapping to a single HP CPU interrupt (i.e., shared interrupts)
    \item Delegation of CPU User Mode interrupts to Machine Mode interrupts
\end{itemize}

\section{Architecture}

Figure \ref{fig:interrupt-matrix-structure} shows the structure of the interrupt matrix.

You need to configure the \hyperref[sec:intmtrx-reg-summ]{interrupt matrix registers} to map the peripheral interrupt sources to the HP CPU interrupts. The Interrupt Matrix Controller in Figure \ref{fig:interrupt-matrix-structure} manages the mapping and sends the interrupt status of each interrupt source to the interrupt status registers which belong to the interrupt matrix registers.

\begin{figure}[H]
    \centering
    \includegraphics[width=0.8\textwidth]{\modulefiles/figures/ESP32-H2-interrupt.eps}
    \caption{Interrupt Matrix Structure}
    \label{fig:interrupt-matrix-structure}
\end{figure}

\section{Functional Description}

\subsection{Peripheral Interrupt Sources}

The \chipname{} has 84 peripheral interrupt sources in total. Table \ref{tab:peripheral-interrupt-config-status-source} lists all these sources and their mapping registers, as well as the interrupt status registers.

\begin{itemize}
    \item Column ``No.": Peripheral interrupt source number, can be 0 \verb+~+ 83.
    \item Column ``Chapter": in which chapter the interrupt source is described in detail.
    \item Column ``Interrupt Source": Name of the peripheral interrupt source.
    \item Column ``Interrupt Source Mapping Register": Registers used to configure the routing of the peripheral interrupt sources to the HP CPU peripheral interrupts.
    \item Column ``Interrupt Status Register": Registers used to reflect the interrupt source status.
    \begin{itemize}
        \item Column ``Interrupt Status Register - Bit": Bit position in status registers, indicating the interrupt status.
        \item Column ``Interrupt Status Register - Name": Name of status registers.
    \end{itemize}
\end{itemize}

\clearpage
\newgeometry{top=2cm,bottom=2cm,left=1.5cm,right=1.5cm}

\begin{landscape}
\begin{footnotesize}
\begin{longtable}[c]{|*{6}{l|C{5cm}|C{3.7cm}|C{7.8cm}|p{0.3cm}|C{4.8cm}|}}
    \caption{CPU Peripheral Interrupt Source Mapping/Status Registers and Peripheral Interrupt Sources}
    \label{tab:peripheral-interrupt-config-status-source}\\

    \hline
    \rowcolor{gray!25}
    & & & & \multicolumn{2}{c|}{\textbf{Interrupt Status Register}} \\
    \cline{5-6}
    \rowcolor{gray!25}
    \multirow{-2}*{\textbf{No.}} & \multirow{-2}*{\textbf{Chapter}} & \multirow{-2}*{\textbf{Interrupt Source}} & \multirow{-2}*{\textbf{Interrupt Source Mapping Register}} & \textbf{Bit} & \textbf{Name} \\
    \hline
    \endfirsthead

    \hline
    \rowcolor{gray!25}
    & & & & \multicolumn{2}{c|}{\textbf{Interrupt Status Register}} \\
    \cline{5-6}
    \rowcolor{gray!25}
    \multirow{-2}*{\textbf{No.}} & \multirow{-2}*{\textbf{Chapter}} & \multirow{-2}*{\textbf{Interrupt Source}} & \multirow{-2}*{\textbf{Interrupt Source Mapping Register}} & \textbf{Bit} & \textbf{Name} \\
    \hline
    \endhead
    \endfoot

    0 &  n/a & reserved & reserved & 0 & \multirow{32}*{\makecell{\hyperref[regdesc:INTMTXCORE0INTSTATUS0REG]{INTMTX\_CORE0\_INT\_STATUS\_0\_REG}}}  \\
    \cline{1-5}
    1 &  n/a & reserved & reserved & 1 & \\
    \cline{1-5}
    2 &  n/a & reserved & reserved & 2 & \\
    \cline{1-5}
    3 &  n/a & reserved & reserved & 3 & \\
    \cline{1-5}
    4 &  n/a & reserved & reserved & 4 & \\
    \cline{1-5}
    5 &  n/a & reserved & reserved & 5 & \\
    \cline{1-5}
    6 &  n/a & reserved & reserved & 6 & \\
    \cline{1-5}
    7 & n/a & reserved & reserved & 7 & \\
    \cline{1-5}
    8 &  n/a & reserved & reserved & 9 & \\
    \cline{1-5}
    9 &  n/a & reserved & reserved & 9 & \\
    \cline{1-5}
    10 &  n/a & reserved & reserved & 10 & \\
    \cline{1-5}
    11 &  n/a & reserved & reserved & 11 & \\
    \cline{1-5}
    12 &  n/a & reserved & reserved & 12 & \\
    \cline{1-5}
    13 & \textit{\nameref{mod:lowpowm}} & PMU\_INTR &\hyperref[regdesc:INTMTXCORE0PMUINTRMAPREG]{INTMTX\_CORE0\_PMU\_INTR\_MAP\_REG} & 13 &                        	\\
    \cline{1-5}
    14  & \textit{\nameref{mod:efuse}} & EFUSE\_INTR  & \hyperref[regdesc:INTMTXCORE0EFUSEINTRMAPREG]{INTMTX\_CORE0\_EFUSE\_INTR\_MAP\_REG} & 14 &                          	\\
    \cline{1-5}
    15 & \textit{\nameref{mod:lowpowm}} & LP\_RTC\_TIMER\_INTR &\hyperref[regdesc:INTMTXCORE0LPRTCTIMERINTRMAPREG]{INTMTX\_CORE0\_LP\_RTC\_TIMER\_INTR\_MAP\_REG} & 15 & \\
    \cline{1-5}
    16 & \textit{\nameref{mod:uart}} & LP\_UART\_INTR &\hyperref[regdesc:INTMTXCORE0LPUARTINTRMAPREG]{INTMTX\_CORE0\_LP\_UART\_INTR\_MAP\_REG} & 16 & \\
    \cline{1-5}
    17 & \textit{\nameref{mod:i2c}} & LP\_I2C\_INTR &\hyperref[regdesc:INTMTXCORE0LPI2CINTRMAPREG]{INTMTX\_CORE0\_LP\_I2C\_INTR\_MAP\_REG} & 17 & \\
    \cline{1-5}
    18 & \textit{\nameref{mod:lowpowm}} & LP\_WDT\_INTR &\hyperref[regdesc:INTMTXCORE0LPWDTINTRMAPREG]{INTMTX\_CORE0\_LP\_WDT\_INTR\_MAP\_REG} & 18 & \\
    \cline{1-5}
    19 & \textit{\nameref{mod:sysreg}} & LP\_PERI\_TIMEOUT\_INTR & \hyperref[regdesc:INTMTXCORE0LPPERITIMEOUTINTRMAPREG]{INTMTX\_CORE0\_LP\_PERI\_TIMEOUT\_INTR\_MAP\_REG} & 19 & \\
    \cline{1-5}
    20 & \textit{\nameref{mod:permctrl}} & LP\_APM\_M0\_INTR & \hyperref[regdesc:INTMTXCORE0LPAPMM0INTRMAPREG]{INTMTX\_CORE0\_LP\_APM\_M0\_INTR\_MAP\_REG} & 20 &                          	\\
    \cline{1-5}
    21 & \textit{\nameref{mod:permctrl}} & LP\_APM\_M1\_INTR & \hyperref[regdesc:INTMTXCORE0LPAPMM1INTRMAPREG]{INTMTX\_CORE0\_LP\_APM\_M1\_INTR\_MAP\_REG} & 21 &                          	\\
    \cline{1-5}
    22 &  \textit{\nameref{mod:keymng}} & HUK\_INTR & \hyperref[regdesc:INTMTXCORE0HUKINTRMAPREG]{INTMTX\_CORE0\_HUK\_INTR\_MAP\_REG} & 22 & \\
    \cline{1-5}
    23 & \textit{\nameref{sec:intpri-reg}} & CPU\_INTR\_FROM\_CPU\_0\label{CPUINTRFROMCPU0} & \hyperref[regdesc:INTMTXCORE0CPUINTRFROMCPU0MAPREG]{INTMTX\_CORE0\_CPU\_INTR\_FROM\_CPU\_0\_MAP\_REG} & 23 &                          	\\
    \cline{1-5}
    24 & \textit{\nameref{sec:intpri-reg}} & CPU\_INTR\_FROM\_CPU\_1 & \hyperref[regdesc:INTMTXCORE0CPUINTRFROMCPU1MAPREG]{INTMTX\_CORE0\_CPU\_INTR\_FROM\_CPU\_1\_MAP\_REG} & 24 &                          	\\
    \cline{1-5}
    25 & \textit{\nameref{sec:intpri-reg}} & CPU\_INTR\_FROM\_CPU\_2 & \hyperref[regdesc:INTMTXCORE0CPUINTRFROMCPU2MAPREG]{INTMTX\_CORE0\_CPU\_INTR\_FROM\_CPU\_2\_MAP\_REG} & 25 &                          	\\
    \cline{1-5}
    26 & \textit{\nameref{mod:riscvcpu}} & CPU\_INTR\_FROM\_CPU\_3 & \hyperref[regdesc:INTMTXCORE0CPUINTRFROMCPU3MAPREG]{INTMTX\_CORE0\_CPU\_INTR\_FROM\_CPU\_3\_MAP\_REG} & 26 &                          	\\
    \cline{1-5}
    27 & \textit{\nameref{mod:dbgassist}} & BUS\_MONITOR\_INTR  &\hyperref[regdesc:INTMTXCORE0BUSMONITORINTRMAPREG]{INTMTX\_CORE0\_BUS\_MONITOR\_INTR\_MAP\_REG}  & 27  &                          	\\
    \cline{1-5}
    28 & \textit{\nameref{mod:riscvcpu}} & TRACE\_INTR & \hyperref[regdesc:INTMTXCORE0TRACEINTRMAPREG]{INTMTX\_CORE0\_TRACE\_INTR\_MAP\_REG} & 28 &                          	\\
    \cline{1-5}
    29 &  n/a & reserved & reserved & 29 & \\
    \cline{1-5}
    30 & \textit{\nameref{mod:sysreg}} & CPU\_PERI\_TIMEOUT\_INTR  & \hyperref[regdesc:INTMTXCORE0CPUPERITIMEOUTINTRMAPREG]{INTMTX\_CORE0\_CPU\_PERI\_TIMEOUT\_INTR\_MAP\_REG}  & 30 &                          	\\
    \cline{1-5}
    31 & \textit{\nameref{mod:iomuxgpio}} & GPIO\_INTR\_PRO & \hyperref[regdesc:INTMTXCORE0GPIOINTRPROMAPREG]{INTMTX\_CORE0\_GPIO\_INTR\_PRO\_MAP\_REG} & 31 &                          	\\
    \hline
    32 & \textit{\nameref{mod:iomuxgpio}} & GPIO\_INTR\_EXT & \hyperref[regdesc:INTMTXCORE0GPIOINTREXTMAPREG]{INTMTX\_CORE0\_GPIO\_INTR\_EXT\_MAP\_REG} & 0 & \hyperref[regdesc:INTMTXCORE0INTSTATUS1REG]{INTMTX\_CORE0\_INT\_STATUS\_1\_REG}\\
    \cline{1-5}
    33 & \textit{\nameref{mod:lowpowm}} & PAU\_INTR &\hyperref[regdesc:INTMTXCORE0PAUINTRMAPREG]{INTMTX\_CORE0\_PAU\_INTR\_MAP\_REG} & 1 & \\
    \cline{1-5}
    34 & \textit{\nameref{mod:sysreg}} & HP\_PERI\_TIMEOUT\_INTR  & \hyperref[regdesc:INTMTXCORE0CPUINTRFROMCPU0MAPREG]{INTMTX\_CORE0\_HP\_PERI\_TIMEOUT\_INTR\_MAP\_REG}  & 2 &                          	\\
    \cline{1-5}
    35 & n/a & reserved & reserved & 3 & \\
    \cline{1-5}
    36 & \textit{\nameref{mod:permctrl}} & HP\_APM\_M0\_INTR & \hyperref[regdesc:INTMTXCORE0HPAPMM0INTRMAPREG]{INTMTX\_CORE0\_HP\_APM\_M0\_INTR\_MAP\_REG} & 4 &                          	\\
    \cline{1-5}
    37 & \textit{\nameref{mod:permctrl}} & HP\_APM\_M1\_INTR & \hyperref[regdesc:INTMTXCORE0HPAPMM1INTRMAPREG]{INTMTX\_CORE0\_HP\_APM\_M1\_INTR\_MAP\_REG} & 5 &                          	\\
    \cline{1-5}
    38 & \textit{\nameref{mod:permctrl}} & HP\_APM\_M2\_INTR & \hyperref[regdesc:INTMTXCORE0HPAPMM2INTRMAPREG]{INTMTX\_CORE0\_HP\_APM\_M2\_INTR\_MAP\_REG} & 6 &                          	\\
    \cline{1-5}
    39 & \textit{\nameref{mod:permctrl}} & HP\_APM\_M3\_INTR & \hyperref[regdesc:INTMTXCORE0HPAPMM3INTRMAPREG]{INTMTX\_CORE0\_HP\_APM\_M3\_INTR\_MAP\_REG} & 7 &                          	\\
    \cline{1-5}
    40 & \textit{\nameref{mod:permctrl}} & HP\_APM\_M4\_INTR & \hyperref[regdesc:INTMTXCORE0HPAPMM4INTRMAPREG]{INTMTX\_CORE0\_HP\_APM\_M4\_INTR\_MAP\_REG} & 8 &                          	\\
    \cline{1-5}
    41 & \textit{\nameref{mod:permctrl}} & LP\_APM0\_INTR & \hyperref[regdesc:INTMTXCORE0LPAPM0INTRMAPREG]{INTMTX\_CORE0\_LP\_APM0\_INTR\_MAP\_REG} & 9 &                          	\\
    \cline{1-5}
    42 & \textit{\nameref{mod:permctrl}} & CPU\_APM\_M0\_INTR & \hyperref[regdesc:INTMTXCORE0CPUAPMM0INTRMAPREG]{INTMTX\_CORE0\_CPU\_APM\_M0\_INTR\_MAP\_REG} & 10 &                          	\\
    \cline{1-5}
    43 & \textit{\nameref{mod:permctrl}} & CPU\_APM\_M1\_INTR & \hyperref[regdesc:INTMTXCORE0CPUAPMM1INTRMAPREG]{INTMTX\_CORE0\_CPU\_APM\_M1\_INTR\_MAP\_REG} & 11 &                          	\\
    \cline{1-5}
    44 & \textit{\nameref{mod:permctrl}} & MSPI\_INTR & \hyperref[regdesc:INTMTXCORE0MSPIINTRMAPREG]{INTMTX\_CORE0\_MSPI\_INTR\_MAP\_REG} & 12 &                          	\\
    \cline{1-5}
    45 & \textit{\nameref{mod:i2s}} & I2S\_INTR & \hyperref[regdesc:INTMTXCORE0I2SINTRMAPREG]{INTMTX\_CORE0\_I2S\_INTR\_MAP\_REG} & 13 &                          	\\
    \cline{1-5}

    46 & \textit{\nameref{mod:uart}} & UHCI0\_INTR & \hyperref[regdesc:INTMTXCORE0UHCI0INTRMAPREG]{INTMTX\_CORE0\_UHCI0\_INTR\_MAP\_REG} & 14 &                           	\\
    \cline{1-5}
    47 & \textit{\nameref{mod:uart}} & UART0\_INTR & \hyperref[regdesc:INTMTXCORE0UART0INTRMAPREG]{INTMTX\_CORE0\_UART0\_INTR\_MAP\_REG} & 15 &                          	\\
    \cline{1-5}
    48 & \textit{\nameref{mod:uart}} & UART1\_INTR & \hyperref[regdesc:INTMTXCORE0UART1INTRMAPREG]{INTMTX\_CORE0\_UART1\_INTR\_MAP\_REG} & 16 &                          	\\
    \cline{1-5}
    49  & \textit{\nameref{mod:ledpwm}} & LEDC\_INTR  & \hyperref[regdesc:INTMTXCORE0LEDCINTRMAPREG]{INTMTX\_CORE0\_LEDC\_INTR\_MAP\_REG} & 17 &                          	\\
    \cline{1-5}
    50  & \textit{\nameref{mod:twai}} & TWAI0\_TIMER\_INTR & \hyperref[regdesc:INTMTXCORE0TWAI0TIMERINTRMAPREG]{INTMTX\_CORE0\_TWAI0\_TIMER\_INTR\_MAP\_REG} & 18 &                          	\\
    \cline{1-5}
    51  & \textit{\nameref{mod:twai}} & TWAI0\_INTR & \hyperref[regdesc:INTMTXCORE0TWAI0INTRMAPREG]{INTMTX\_CORE0\_TWAI0\_INTR\_MAP\_REG} & 19 &                          	\\
    \cline{1-5}
    52  & \textit{\nameref{mod:twai}} & TWAI1\_TIMER\_INTR & \hyperref[regdesc:INTMTXCORE0TWAI1TIMERINTRMAPREG]{INTMTX\_CORE0\_TWAI1\_TIMER\_INTR\_MAP\_REG} & 20&                          	\\
    \cline{1-5}
    53  & \textit{\nameref{mod:twai}} & TWAI1\_INTR & \hyperref[regdesc:INTMTXCORE0TWAI1INTRMAPREG]{INTMTX\_CORE0\_TWAI1\_INTR\_MAP\_REG} & 21 &                          	\\
    \cline{1-5}
    54  & \textit{\nameref{mod:usbserialjtag}} & USB\_SERIAL\_JTAG\_INTR  & \hyperref[regdesc:INTMTXCORE0USBSERIALJTAGINTRMAPREG]{INTMTX\_CORE0\_USB\_INTR\_MAP\_REG}  & 22 &                          	\\
    \cline{1-5}
    55  & \textit{\nameref{mod:rmt}} &  RMT\_INTR & \hyperref[regdesc:INTMTXCORE0RMTINTRMAPREG]{INTMTX\_CORE0\_RMT\_INTR\_MAP\_REG} & 23  &                          	\\
    \cline{1-5}
    56  & \textit{\nameref{mod:i2c}} & I2C\_EXT0\_INTR  & \hyperref[regdesc:INTMTXCORE0I2CEXT0INTRMAPREG]{INTMTX\_CORE0\_I2C\_EXT0\_INTR\_MAP\_REG}  & 24  &                          	\\
    \cline{1-5}
    57 & \textit{\nameref{mod:timg}} & TG0\_T0\_INTR & \hyperref[regdesc:INTMTXCORE0TG0T0INTRMAPREG]{INTMTX\_CORE0\_TG0\_T0\_INTR\_MAP\_REG} & 25 &                      	\\
    \cline{1-5}
    58  & \textit{\nameref{mod:timg}} & TG0\_WDT\_INTR & \hyperref[regdesc:INTMTXCORE0TG0WDTINTRMAPREG]{INTMTX\_CORE0\_TG0\_WDT\_INTR\_MAP\_REG}  & 26 &                          	\\
    \cline{1-5}
    59  & \textit{\nameref{mod:timg}} & TG1\_T0\_INTR & \hyperref[regdesc:INTMTXCORE0TG1T0INTRMAPREG]{INTMTX\_CORE0\_TG1\_T0\_INTR\_MAP\_REG}  & 27 &                          	\\
    \cline{1-5}
    60  & \textit{\nameref{mod:timg}} & TG1\_WDT\_INTR & \hyperref[regdesc:INTMTXCORE0TG1WDTINTRMAPREG]{INTMTX\_CORE0\_TG1\_WDT\_INTR\_MAP\_REG}  & 28 &                          	\\
    \cline{1-5}
    61 & \textit{\nameref{mod:systimer}}  & SYSTIMER\_TARGET0\_INTR & \hyperref[regdesc:INTMTXCORE0SYSTIMERTARGET0INTRMAPREG]{INTMTX\_CORE0\_SYSTIMER\_TARGET0\_INTR\_MAP\_REG}  & 29 &                          	\\
    \cline{1-5}
    62 & \textit{\nameref{mod:systimer}} & SYSTIMER\_TARGET1\_INTR &\hyperref[regdesc:INTMTXCORE0SYSTIMERTARGET1INTRMAPREG]{INTMTX\_CORE0\_SYSTIMER\_TARGET1\_INTR\_MAP\_REG}  & 30 &                            	\\
    \cline{1-5}
    63  & \textit{\nameref{mod:systimer}} & SYSTIMER\_TARGET2\_INTR &\hyperref[regdesc:INTMTXCORE0SYSTIMERTARGET2INTRMAPREG]{INTMTX\_CORE0\_SYSTIMER\_TARGET2\_INTR\_MAP\_REG}  & 31 &                          	\\
    \hline
    64 & \textit{\nameref{mod:adc}} & APB\_ADC\_INTR & \hyperref[regdesc:INTMTXCORE0APBADCINTRMAPREG]{INTMTX\_CORE0\_APB\_ADC\_INTR\_MAP\_REG}  & 0  & \hyperref[regdesc:INTMTXCORE0INTSTATUS2REG]{INTMTX\_CORE0\_INT\_STATUS\_2\_REG}\\
    \cline{1-5}
    65 & \textit{\nameref{mod:mcpwm}}  & PWM\_INTR  & \hyperref[regdesc:INTMTXCORE0PWMINTRMAPREG]{INTMTX\_CORE0\_PWM\_INTR\_MAP\_REG} & 1  & \\
    \cline{1-5}
    66 & \textit{\nameref{mod:pcnt}}  & PCNT\_INTR  & \hyperref[regdesc:INTMTXCORE0PCNTINTRMAPREG]{INTMTX\_CORE0\_PCNT\_INTR\_MAP\_REG} & 2  & \\
    \cline{1-5}
    67 & \textit{\nameref{mod:parlio}}  & PARL\_IO\_TX\_INTR  & \hyperref[regdesc:INTMTXCORE0PARLIOTXINTRMAPREG]{INTMTX\_CORE0\_PARL\_IO\_TX\_INTR\_MAP\_REG} & 3  & \\
    \cline{1-5}
    68 & \textit{\nameref{mod:parlio}}  & PARL\_IO\_RX\_INTR  & \hyperref[regdesc:INTMTXCORE0PARLIORXINTRMAPREG]{INTMTX\_CORE0\_PARL\_IO\_RX\_INTR\_MAP\_REG} & 4  & \\\cline{1-5}
    69 & \textit{\nameref{mod:sdioslave}} & SLC0\_INTR  & \hyperref[regdesc:INTMTXCORE0SLC0INTRMAPREG]{INTMTX\_CORE0\_SLC0\_INTR\_MAP\_REG}  & 5 &                          	\\
    \cline{1-5}
    70 & \textit{\nameref{mod:sdioslave}} & SLC1\_INTR  & \hyperref[regdesc:INTMTXCORE0SLC1INTRMAPREG]{INTMTX\_CORE0\_SLC1\_INTR\_MAP\_REG}  & 6 &                          	\\
    \cline{1-5}
    71 & \textit{\nameref{mod:dma}} & GDMA\_IN\_CH0\_INTR  & \hyperref[regdesc:INTMTXCORE0DMAINCH0INTRMAPREG]{INTMTX\_CORE0\_DMA\_IN\_CH0\_INTR\_MAP\_REG}  & 7 &                          	\\
    \cline{1-5}
    72 & \textit{\nameref{mod:dma}} & GDMA\_IN\_CH1\_INTR  & \hyperref[regdesc:INTMTXCORE0DMAINCH1INTRMAPREG]{INTMTX\_CORE0\_DMA\_IN\_CH1\_INTR\_MAP\_REG}  & 8 &                          	\\
    \cline{1-5}
    73 & \textit{\nameref{mod:dma}} & GDMA\_IN\_CH2\_INTR  & \hyperref[regdesc:INTMTXCORE0DMAINCH2INTRMAPREG]{INTMTX\_CORE0\_DMA\_IN\_CH2\_INTR\_MAP\_REG}  & 9 &                          	\\
    \cline{1-5}
    74 & \textit{\nameref{mod:dma}} & GDMA\_OUT\_CH0\_INTR  & \hyperref[regdesc:INTMTXCORE0DMAOUTCH0INTRMAPREG]{INTMTX\_CORE0\_DMA\_OUT\_CH0\_INTR\_MAP\_REG}  & 10 &                          	\\
    \cline{1-5}
    75 & \textit{\nameref{mod:dma}} & GDMA\_OUT\_CH1\_INTR  & \hyperref[regdesc:INTMTXCORE0DMAOUTCH1INTRMAPREG]{INTMTX\_CORE0\_DMA\_OUT\_CH1\_INTR\_MAP\_REG}  & 11 &                          	\\
    \cline{1-5}
    76 & \textit{\nameref{mod:dma}} & GDMA\_OUT\_CH2\_INTR  & \hyperref[regdesc:INTMTXCORE0DMAOUTCH2INTRMAPREG]{INTMTX\_CORE0\_DMA\_OUT\_CH2\_INTR\_MAP\_REG}  & 12 &                          	\\
    \cline{1-5}
    77 & \textit{\nameref{mod:spi}} & GPSPI2\_INTR & \hyperref[regdesc:INTMTXCORE0GPSPI2INTRMAPREG]{INTMTX\_CORE0\_GPSPI2\_INTR\_MAP\_REG} & 13 &                          	\\
    \cline{1-5}
    78 & \textit{\nameref{mod:aes}} & AES\_INTR  & \hyperref[regdesc:INTMTXCORE0AESINTRMAPREG]{INTMTX\_CORE0\_AES\_INTR\_MAP\_REG}  & 14 &                          	\\
    \cline{1-5}
    79 & \textit{\nameref{mod:sha}} & SHA\_INTR  & \hyperref[regdesc:INTMTXCORE0SHAINTRMAPREG]{INTMTX\_CORE0\_SHA\_INTR\_MAP\_REG}  & 15 &                          	\\
    \cline{1-5}
    80 & \textit{\nameref{mod:rsa}} & RSA\_INTR  & \hyperref[regdesc:INTMTXCORE0RSAINTRMAPREG]{INTMTX\_CORE0\_RSA\_INTR\_MAP\_REG}  & 16 &                          	\\
    \cline{1-5}
    81 & \textit{\nameref{mod:ecc}} & ECC\_INTR  & \hyperref[regdesc:INTMTXCORE0ECCINTRMAPREG]{INTMTX\_CORE0\_ECC\_INTR\_MAP\_REG}  & 17 &                          	\\
    \cline{1-5}
    82 & \textit{\nameref{mod:ecdsa}} & ECDSA\_INTR  & \hyperref[regdesc:INTMTXCORE0ECDSAINTRMAPREG]{INTMTX\_CORE0\_ECDSA\_INTR\_MAP\_REG}  & 18 &\\
    \cline{1-5}
    83 & \textit{\nameref{mod:keymng}} & KM\_INTR & \hyperref[regdesc:INTMTXCORE0KMINTRMAPREG]{INTMTX\_CORE0\_KM\_INTR\_MAP\_REG}  & 19 &                          	\\
    \hline

\end{longtable}
\end{footnotesize}

\end{landscape}
\restoregeometry
\clearpage

\subsection{HP CPU Interrupts}

The HP CPU of the \chipname{} implements its interrupt mechanism using standard RISC-V (v0.9) Core-Local Interrupt Controller (CLIC) interrupt. The HP CPU supports 2 core local interrupts (CLINT) and 32 peripheral interrupts, which are numbered from 16 to 47. The peripheral interrupt sources are mapped to the 32 HP CPU peripheral interrupts through the interrupt matrix. Each HP CPU peripheral interrupt has the following properties:

\begin{itemize}
    \item Priority levels configurable, see Chapter \ref{mod:riscvcpu} \textit{\nameref{mod:riscvcpu}} > Section \ref{sec:riscvcpu-clic-int-level-prio} \textit{\nameref{sec:riscvcpu-clic-int-level-prio}} for details. 
    \item Configurable as level-triggered or edge-triggered.
    \item Lower-priority interrupts maskable by setting interrupt threshold.
\end{itemize}

\begin{tiplisting}
For detailed information about the functions and configuration procedures of CPU interrupts, see Chapter \ref{mod:riscvcpu} \textit{\nameref{mod:riscvcpu}} > Section \ref{sec:riscvcpu-int-ctrl}\textit{\nameref{sec:riscvcpu-int-ctrl}}. The configuration registers of CPU interrupts are listed in Section \ref{sec:intpri-reg-summ} \textit{\nameref{sec:intpri-reg-summ}}.
\end{tiplisting}
\vspace{-1em}

\subsection{Assign Peripheral Interrupt Source to HP CPU Peripheral Interrupt}

In this section, the following terms are used to describe the operation of the interrupt matrix.

\begin{itemize}

\item \regindex{SOURCE}: stands for a peripheral interrupt source in Table \ref{tab:peripheral-interrupt-config-status-source}.

\item \hyperref[INTMTXCORE0SOURCEINTRMAPREG]{INTMTX\_CORE0\_\regindex{SOURCE}\_INTR\_MAP\_REG}: stands for the interrupt source mapping register for a \regindex{SOURCE}.

\item Num\_P: stands for the index of HP CPU interrupts which can be 16 \verb+~+ 47.

\item Interrupt\_P: stands for the HP CPU interrupt numbered as Num\_P.
\end{itemize}

\subsubsection{Assign One Peripheral Interrupt Source to HP CPU Peripheral Interrupt}

Setting the corresponding source mapping register \hyperref[INTMTXCORE0SOURCEINTRMAPREG]{INTMTX\_CORE0\_\regindex{SOURCE}\_INTR\_MAP\_REG} of \regindex{SOURCE} to Num\_P assigns this interrupt source to Interrupt\_P.

\subsubsection{Assign Multiple Peripheral Interrupt Sources to HP CPU Peripheral Interrupt}

Setting the corresponding source mapping register \hyperref[INTMTXCORE0SOURCEINTRMAPREG]{INTMTX\_CORE0\_\regindex{SOURCE}\_INTR\_MAP\_REG} of each \regindex{SOURCE} to the same Num\_P assigns multiple sources to the same Interrupt\_P. Any of these sources can trigger CPU Interrupt\_P. When an interrupt signal is generated, the CPU should check the interrupt status registers to figure out which peripheral generated the interrupt. For more information, see Chapter \ref{mod:riscvcpu} \textit{\nameref{mod:riscvcpu}} > Section \ref{sec:riscvcpu-int-ctrl}\textit{\nameref{sec:riscvcpu-int-ctrl}}.

\subsubsection{Unassign SOURCE}

Writing 0 to the \hyperref[INTMTXCORE0SOURCEINTRMAPREG]{INTMTX\_CORE0\_\regindex{SOURCE}\_INTR\_MAP\_REG} register disables the corresponding interrupt source.


\subsection{Delegated Interrupts}

The ESP32-C5 RISC-V CPU supports Machine Mode and User Mode. The 32 external HP CPU interrupts can be configured as either Machine Mode interrupts or User Mode interrupts. When the HP CPU is in Machine Mode, it only responds to Machine Mode interrupts and does not respond to User Mode interrupts. When the HP CPU is in User Mode, it can respond to all interrupts.

Under normal conditions, User Mode interrupts are only handled when the CPU is operating in User Mode. However, with the interrupt delegation feature, Machine Mode can be configured to handle certain User Mode interrupts, while the CPU itself remains in Machine Mode. This is achieved by configuring registers to delegate User Mode interrupts to Machine Mode interrupts.

\begin{itemize}
\item Each interrupt source of the interrupt matrix can independently enable or disable the interrupt delegation feature.
\item Multiple interrupt sources from the interrupt matrix can enable the interrupt delegation feature at the same time.
\item Delegation is allowed to only one interrupt signal from the interrupt matrix.
\end{itemize}

Software can configure the \hyperref[regdesc:INTMTXCORE0INTSIGIDXASSERTINSECREG]{INTMTX\_CORE0\_INT\_SIG\_IDX\_ASSERT\_IN\_SEC\_REG} register to n to delegate a User Mode interrupt to a Machine Mode interrupt with interrupt number n.

To enable the delegated interrupt function for interrupt source \regindex{SOURCE} in the interrupt matrix, software can set the corresponding bit in the \hyperref[INTMTXCORE0SOURCEINTRMAPREG]{INTMTX\_CORE0\_\regindex{SOURCE}\_INTR\_PASS\_IN\_SEC\_REG} register to 1.

\subsection{Query Current Interrupt Status of SOURCE}

After enabling a \regindex{SOURCE}, you can query its current interrupt status by reading the corresponding bit value in \hyperref[regdesc:INTMTXCORE0INTSTATUS0REG]{INTMTX\_CORE0\_INT\_STATUS\_\regindex{n}\_REG} (read only). For the mapping between \hyperref[regdesc:INTMTXCORE0INTSTATUS0REG]{INTMTX\_CORE0\_INT\_STATUS\_\regindex{n}\_REG} and the \regindex{SOURCE}, please refer to Table \ref{tab:peripheral-interrupt-config-status-source}.


% >>> If your module has no registers, delete sample text
%       for Sections Register Summary and Registers below <<<









%%%%%%%%%%%%%%%%%%%%%%%%
%%% Register Summary %%%
%%%%%%%%%%%%%%%%%%%%%%%%

% >>> Update the sample text below and insert
%     the info on register summary into the subfile <<<

\newpage
\hypertarget{intmtrx-reg-summ}{}
\begin{landscape}
\section{Register Summary}


\subsection{Interrupt Matrix Register Summary}
\label{sec:intmtrx-reg-summ}

The addresses in this section are relative to the interrupt matrix base address provided in Table \ref{tab:sysmem-base-address} in Chapter \ref{mod:sysmem} \textit{\nameref{mod:sysmem}}.

The abbreviations given in Column \textbf{Access} are explained in Section \hyperref[glossary-access-types]{\textit{Access Types for Registers}}.

\subfile{\modulefiles/11-interrupt_matrix_reg_summary__EN}


\subsection{Software Interrupt Register Summary}
\label{sec:intpri-reg-summ}
The addresses in this section are relative to the software interrupt base address provided in Table \ref{tab:sysmem-base-address} in Chapter \ref{mod:sysmem} \textit{\nameref{mod:sysmem}}.

The abbreviations given in Column \textbf{Access} are explained in Section \hyperref[glossary-access-types]{\textit{Access Types for Registers}}.

\subfile{\modulefiles/11-intpri_reg_summary__EN}

\end{landscape}

%%%%%%%%%%%%%%%%%
%%% Registers %%%
%%%%%%%%%%%%%%%%%

% >>> Update the sample text below and insert
%      the info on registers into the subfile <<<

\newpage
\section{Registers}

\subsection{Interrupt Matrix Registers}

The addresses in this section are relative to  the interrupt matrix base address provided in Table \ref{tab:sysmem-base-address} in Chapter \ref{mod:sysmem} \textit{\nameref{mod:sysmem}}.

\subfile{\modulefiles/11-interrupt_matrix_reg__EN}

\subsection{Software Interrupt Registers}
\label{sec:intpri-reg}

The addresses in this section are relative to  the software interrupt base address provided in Table \ref{tab:sysmem-base-address} in Chapter \ref{mod:sysmem} \textit{\nameref{mod:sysmem}}.

\subfile{\modulefiles/11-intpri_reg__EN}


\end{document}
